\paper{非数学类}

\begin{problems}
\item[一、]（本题 30 分，每小题 6 分）填空题：
\begin{enumerate}
  \item 若 $f(x)$ 在点 $x=a$ 处可导，且 $f(a)\neq 0$，则 $\displaystyle \lim_{n\to+\infty}\left[\frac{f(a+1/n)}{f(a)}\right]^{n}=\underline{\hspace{4em}}$．
  \item 若 $f(1)=0$，$f'(1)$ 存在，则极限 $\displaystyle I=\lim_{x\to 0}\frac{f(\sin^2x+\cos x)\tan 3x}{\left(\ee^{x^2}-1\right)\sin x}=\underline{\hspace{4em}}$．
  \item 设 $f(x)$ 有连续导数，且 $f(1)=2$．记 $z=f(\ee^xy^2)$，若 $\displaystyle \frac{\partial z}{\partial x}=z$，$f(x)$ 在 $x>0$ 的表达式为 $\underline{\hspace{6em}}$．
  \item 设 $f(x)=\ee^x\sin 2x$，则 $f^{(4)}(0)=\underline{\hspace{6em}}$．
  \item 曲面 $z=\frac{x^2}{2}+y^2$ 平行于平面 $2x+2y-z=0$ 的切平面方程为 $\underline{\hspace{8em}}$．
\end{enumerate}

\item[二、]（本题 14 分）设 $f(x)$ 在 $[0,1]$ 上可导，$f(0)=0$，且当 $x\in(0,1)$，$0<f'(x)<1$．试证：当 $a\in(0,1)$ 时，有
\[
  \left(\int_0^a f(x)\dd x\right)^2>\int_0^a f^3(x)\dd x.
\]

\item[三、]（本题 14 分）某物体所在的空间区域为 $\Omega$：$x^2+y^2+2z^2\le x+y+2z$，密度函数为 $x^2+y^2+z^2$，求质量
\[
  M=\iiint_\Omega \left(x^2+y^2+z^2\right)\dd x\dd y\dd z.
\]

\item[四、]（本题 14 分）设函数 $f(x)$ 在闭区间 $[0,1]$ 上具有连续导数，$f(0)=0$，$f(1)=1$．证明：
\[
  \lim_{n\to\infty}n\left(\int_0^1 f(x)\dd x-\frac{1}{n}\sum_{k=1}^n f\left(\frac{k}{n}\right)\right)=-\frac{1}{2}.
\]

\item[五、]（本题 14 分）设函数 $f(x)$ 在区间 $[0,1]$ 上连续，且 $\displaystyle I=\int_0^1 f(x)\dd x\neq 0$．证明：在 $(0,1)$ 内存在不同的两点 $x_1,x_2$，使得
\[
  \frac{1}{f(x_1)}+\frac{1}{f(x_2)}=\frac{2}{I}.
\]

\item[六、]（本题 14 分）设 $f(x)$ 在 $(-\infty,+\infty)$ 上可导，且 $f(x)=f(x+2)=f(x+\sqrt{3})$，用傅里叶（Fourier）级数理论证明 $f(x)$ 为常数．
\end{problems}
